\section{Introduction}
\label{sec:intro}

Diffusion models have become a dominant paradigm for high-quality image and
multimodal generation, powering modern systems such as Stable Diffusion and
large-scale class-conditional generators.
While most progress has focused on improving training objectives and model
architectures, an emerging line of work shows that \emph{inference-time
scaling}---allocating more computation at test time without retraining---can
sharply improve sample quality under a fixed model.
Recent methods show that test-time techniques such as noise search,
trajectory refinement, and verifier-guided sampling can yield substantial gains
over plain diffusion sampling, when additional compute is available.

Every diffusion-family model---whether a DDPM, a score-based SDE, or a
flow-matching ODE---generates samples by traversing a discrete trajectory
of $T$ timesteps.
Inference-time compute can be invested along this trajectory in two
fundamentally different ways:
(i) \emph{lengthening the trajectory} by increasing $T$ (more denoising
steps), which yields quickly diminishing returns, and
(ii) \emph{searching within each timestep} by evaluating multiple noise
candidates and selecting the one that maximizes a verifier score.
The second axis is the focus of this paper: under a fixed budget of function
evaluations (NFE), how should per-timestep search effort be allocated across
the trajectory?

A key limitation of existing inference-time methods is that they
largely treat the denoising trajectory as a homogeneous sequence of steps.
Most approaches either (i) allocate search computation uniformly across
timesteps, or (ii) hard-code search to specific parts of the trajectory
(e.g., only the initial noise or a fixed middle window).
This overlooks a fundamental empirical fact: \emph{different timesteps exhibit
dramatically different sensitivities to noise perturbations}.
In practice, refining noise at some steps yields large downstream improvements
in sample quality, while spending the same computation at other steps has
negligible effect.
As a result, naive per-step allocation wastes a large fraction of the
test-time compute budget.

In this work, we study \emph{where to search} along the diffusion trajectory.
We formalize inference-time noise refinement as a \emph{noise trajectory search}
problem under a fixed number of function evaluations (NFE), and ask how to
\emph{globally schedule} search computation across timesteps.
Our central observation is that effective inference-time scaling requires
reasoning at two distinct levels: \emph{local} decisions about how to search for
better noise at a given timestep, and \emph{global} decisions about how to
distribute computation across the entire denoising process.
While prior work has explored increasingly sophisticated local search operators,
the problem of global per-step compute allocation has only recently begun
to receive attention, with concurrent work proposing either uniform allocation
with implicit adaptation~\citep{ramesh2025nts} or purely reactive online
reallocation~\citep{kim2025rbf}.

To bridge this gap, we propose a unified two-level framework for noise trajectory
search in diffusion models.
At the local level, we abstract a broad family of existing noise refinement
methods as a generic local search operator applied at a single timestep.
At the global level, we introduce a scheduler that decides \emph{when} and
\emph{how much} local search to perform at each timestep under a total NFE
budget.
Within this framework, we show that standard diffusion sampling, best-of-$N$
selection, $\epsilon$-greedy noise search, and particle- or tree-based methods
are special cases with fixed or heuristic scheduling
policies.

Building on this formulation, we propose \textbf{GAINS} (\textbf{G}lobal \textbf{A}daptive \textbf{I}nference-time \textbf{N}oise \textbf{S}cheduling), a global scheduling algorithm that
combines \emph{offline profiling} with \emph{online control}.
Offline, we analyze verifier sensitivity across timesteps to obtain a coarse
allocation of compute tailored to the model and dataset.
Online, we dynamically adjust per-step search using feedback from each instance,
such as score gains and variance, allowing the sampler to stop early on
unproductive steps and reallocate computation where it is most effective.
This controller operates without access to future timesteps and
respects a strict NFE budget.

We evaluate GAINS on Stable Diffusion, EDM, and ImageNet benchmarks under
fixed compute budgets.
Across a wide range of NFE settings and evaluation metrics, GAINS consistently outperforms uniform allocation and naive
online heuristics.
These results show that \emph{global scheduling with per-step awareness}
is an effective and underexplored dimension of inference-time scaling
for diffusion models.

\paragraph{Contributions.}
Our main contributions are:
(i) a unified two-level view of noise trajectory search that separates local
noise refinement from global timestep scheduling, applicable to both
stochastic (DDPM/SDE) and deterministic (flow/ODE) samplers;
(ii) a novel global scheduling algorithm that integrates \emph{offline timestep
profiling} with \emph{online feedback control}, combining learned difficulty
patterns per step with adaptive early stopping per instance---unlike
purely uniform~\citep{ramesh2025nts}, purely
reactive~\citep{kim2025rbf}, or fixed-phase~\citep{zhang2025tscend}
approaches;
and (iii) extensive experiments on both stochastic (Stable Diffusion, EDM)
and flow-based models, showing consistent gains over existing
inference-time methods under identical NFE budgets.
